\section{Plataforma y separación de baselines}
\subsection{Adopción de kirky-arqui}
Se conserva el procesador desarrollado en el curso: su organización de cinco
etapas, dos lecturas de registros, forwarding y Fetch mixto constituyen una
base comprensible para integrar una operación escalar. No se reemplaza el
core por una plataforma externa. Se congela el commit \code{\CoreShort} y
se exportan sus 143 archivos versionados en \code{baseline/upstream/}.

La base es docente y su ISA es parcial. La auditoría dirigida se presenta en
la sección~\ref{sec:results}; el detalle de módulos, líneas y decisiones está
en \code{docs/CORE\_AUDIT.md}. El éxito de los programas históricos no se
interpreta como certificación de RV32IC completo.

\subsection{Correcciones comunes antes de medir}
Las memorias actuales tienen 64 palabras de 32 bits cada una y lectura
combinacional. Los registros ID/EX se habilitan permanentemente. La lógica de
load-use no distingue todos los campos que realmente son registros fuente.
El proyecto experimental requerirá capacidades de memoria suficientes,
clasificación de operandos, detección de instrucciones excluidas e instrumentación
de retiro. Las variantes multiciclo requerirán además retención y validez.

Estos cambios se aplicarán a una base normalizada común. Si una corrección
reduce pausas espurias, su beneficio no debe atribuirse a XQDot4Z. Se revisará
el subconjunto necesario para los kernels, sin convertir el soporte íntegro
de la ISA en un prerrequisito del estudio aritmético.

\begin{center}
\begin{tikzpicture}
\node[block,text width=57mm] (h) at (0,0) {H0: snapshot histórico\\sin modificaciones};
\node[block,fill=nord8,text width=57mm] (common) at (0,-1.7) {Base normalizada común\\correcciones + memoria + retiro};
\draw[arrow] (h)--(common);
\foreach \variantid/\variantx/\varianttext in {b1/-5.7/B1: software,b2/-1.9/B2: MUL,b3/1.9/B3: dot + corrección,d/5.7/D: XQDot4Z}{
 \node[block,text width=29mm,minimum height=15mm] (\variantid) at (\variantx,-3.7) {\varianttext};
 \draw[arrow] (common.south)--(0,-2.65)-|(\variantid.north);
}
\node[note,text width=145mm] at (0,-5) {Las ramas heredan las mismas correcciones. B2, B3 y D disponen del mismo MUL escalar; B3 y D usan protocolos comparables.};
\end{tikzpicture}
\captionof{figure}{Separación entre procedencia histórica y comparación causal. La base normalizada y sus variantes todavía no están implementadas.}
\end{center}

\begin{tabularx}{\textwidth}{@{}lYY@{}}
\toprule
\textbf{ID} & \textbf{Variante prevista} & \textbf{Comparación interpretada}\\\midrule
H0 & Original exportado. & Estado del proyecto de curso, no efecto puro del ISA.\\
B1 & Base normalizada; multiplicación en software. & Aporte respecto de un core sin multiplicador.\\
B2 & Base normalizada + MUL escalar. & Beneficio del procesamiento empacado/fusión.\\
B3 & Producto punto U4\,$\times$\,S8 y \(P-zS_a\). & Beneficio específico de fusionar la corrección.\\
D & XQDot4Z. & Costo y utilidad de la propuesta.\\\bottomrule
\end{tabularx}\par
MUL por sí sola no constituye la extensión M completa. El estado de soporte
de instrucciones y la revisión de la base se fijarán junto con cada resultado.
